% ED26-SRC-063-2008: complete Fordham manuscript chapter8, printed433--504.
% Final2008 published ED26-SRC-063-0208 and other HHL2 chapters remain distinct.
\chapter{Post-Ockham Latin Logic: Supposition, Consequence, and Obligation}
\label{sec:p13-medieval-latin-inference}
\chaptermark{Post-Ockham Latin Logic}
\index{medieval logic}\index{supposition}\index{consequence!medieval}
\index{obligationes}\index{Dutilh Novaes, Catarina}

Fourteenth-century Latin logic connects semantic analysis, the construction
of arguments, and regulated disputation. This treatment follows Catarina
Dutilh Novaes's chapter on logic after Ockham in the author-hosted Fordham
manuscript of \emph{Mediaeval and Renaissance Logic}
\cite[pp.~433--504]{DutilhNovaes2007FordhamChapter8}. Its source order is
historical setting, semantics, consequences, and obligations.
The chapter complements rather than replaces the volume's separate accounts
of Ockham and Buridan, supposition theory, modality, and insolubilia.

The relations, diagrams, and game reconstructions attributed to Novaes
are modern analyses of historical texts. They are not symbolic inscriptions
found in those medieval texts. Reported original formulations, the modern
chapter's interpretation, and the explicitly specified local mathematical
models below have different authority. A medieval text mentioned in the
chapter is not thereby independently established in its complete original
edition or translation.

\section{The institutions and the traditions}
\label{sec:p13-latin-traditions}

Novaes treats the fourteenth century as an intellectually productive period,
not simply a decline interrupted by the Black Death. Logic occupied the
trivium with grammar and rhetoric and was learned early, but its advanced
literature was not confined to elementary teaching. Buridan's long career
in the arts faculty exemplifies sustained logical investigation.
The chapter's six principal genres are insolubilia, modal propositions,
supposition, proofs of propositions, obligations, and consequences.
\emph{Sophismata} are puzzles addressed with those resources, not a further
uniform theory \cite[pp.~433--435]{DutilhNovaes2007FordhamChapter8}.

The historical classification into British and Parisian traditions is an
explanatory aid, not a separation of independent national schools. Regional
organization at Paris, study abroad, travelling masters, and the foundation
of other universities transmitted and modified both traditions.
British teaching influenced Italy; Parisian and Buridanian approaches
spread into Central Europe. Prague had important links with both Paris
and Oxford. Novaes gives the foundations of Prague in 1348, Vienna in 1365
with re-foundation in 1384, Erfurt's papal authorization in 1379 and opening
in 1392, Heidelberg in 1385, and Cologne in 1388/9 as parts of this account.
These dates describe the source's university history, not dates at which
particular logical doctrines were invented.

\subsection{British development and its transmission}

The antecedents of a specifically British tradition remain disputed.
William of Sherwood, Roger Bacon, Simon of Faversham, and Robert Kilwardby
were important thirteenth-century figures with Parisian connections.
The hypothesis that Oxford preserved terminist methods while Paris became
predominantly Modist, and subsequently exported terminism back to Paris,
is reported as contested rather than settled.

Burley combines new logical resources with realist commitments. His early
consequence and obligation treatises and the shorter and longer
\emph{De puritate} supply starting points for Novaes's technical discussion.
Ockham's \emph{Summa logicae} was influential although his period of
logical writing was relatively short. His semantics and full obligation
account are not the subject of this chapter. Adam Wodeham's significance
here concerns the doctrine of \emph{complexe significabile} and the
transmission and criticism of Ockham, not a new consequence calculus.

The Mertonian or Oxford-Calculator group includes Thomas Bradwardine,
William Heytesbury, Richard Billingham, Roger Swyneshed, and Richard
Kilvington. Roger Swyneshed must not be confused with Richard Swyneshed.
Bradwardine's insolubilia, Heytesbury's divided/composite senses and
\emph{Regulae solvendi sophismata}, Billingham's \emph{Speculum puerorum},
Swyneshed's new obligation procedure, and Kilvington's sophisms make
different contributions. Novaes uses the Black Death in 1349 as a
convenient period divider, not as proof that logical development stopped.
The source places Heytesbury's \emph{Regulae} in 1335 and Kilvington's
relevant sophisms before 1325.

Later British-associated authors include Ralph Strode, Richard Lavenham,
Richard Ferrybridge, John Wyclif, and Paul of Venice.
Lavenham's accessible short treatises exemplify the period's genres;
Strode's \emph{Logica} and Ferrybridge's consequence work had wider
influence but less complete modern editions in the source's scholarly
setting. Wyclif's realism links predication, supposition and truth to
metaphysical commitments. Paul of Venice's Oxford training explains his
placement in this intellectual tradition despite his predominantly Italian
career: his \emph{Logica parva}, placed around 1395, and \emph{Logica magna}
preserve earlier doctrines while making new contributions.
Peter of Candia, Henry Hopton, Robert Fland, Martinus Anglicus,
Johannes Venator, Robert Alington and Richard Brinkley further illustrate
its transmission and diversification
\cite[pp.~437--443]{DutilhNovaes2007FordhamChapter8}.

\subsection{The Parisian and continental development}

Novaes qualifies the apparent early-fourteenth-century Parisian hiatus by
the incomplete study of its surviving material. Modist authors include
Thomas of Erfurt, Radulphus Brito and Siger of Courtrai.
Giraldus Odonis's \emph{Logica}, associated with the early 1320s, provides
evidence of terminist activity that complicates a simple Oxford-import
story. Thomas Maulevelt's influence belongs to this problem of transmission;
variant spellings of his name are not proof of several different authors.

Buridan, Nicholas of Autrecourt, Nicholas of Oresme and Gregory of Rimini
belong to the subsequent generation, but their logical roles differ.
Gregory transmitted Oxford work and defended \emph{complexe significabile}.
Buridan's coherent nominalist system differs materially from Ockham's;
the \emph{via Buridanii} is not merely another name for Ockham's doctrines.
The source connects Albert of Saxony with both influences without making
him simply Buridan's pupil. William Buser, Thomas of Cleves and Marsilius
of Inghen extend the transmission story and introduce semantic refinements
examined below.

John of Holland's Prague career and familiarity with Oxford texts do not
settle where he acquired that familiarity. Peter of Ailly's mental-language
and insolubilia contribution remains separate from this chapter's detailed
analysis. Vincent Ferrer's Spanish and Thomist setting makes his recognition
of collective supposition a historical comparison rather than evidence of
a demonstrated direct line of influence from Paris. Blaise of Parme and
John Dorp illustrate Italian reception and the later consolidation of
Buridanian logic. Dorp's 1393 curricular role and his later printed
commentary must not be confused with the dating of every doctrine in it
\cite[pp.~443--448]{DutilhNovaes2007FordhamChapter8}.

\section{Supposition and the analysis of propositions}
\label{sec:p13-latin-supposition}

\subsection{Personal, simple, and material are author-specific classifications}

Personal supposition concerns things signified by a term; material
supposition concerns terms or term occurrences. In the older realist
classification, simple supposition concerns a corresponding universal.
Ockham rejects the relevant extra-mental universal but preserves simple
supposition as supposition for a mental concept. Buridan instead places
supposition for mental terms together with other material supposition
and allows only personal supposition in mental language.
These are distinct classifications, not one fixed taxonomy.

Albert of Saxony gives a further variant. Spoken and written terms can
have personal, simple and material supposition; mental terms have personal
and material supposition. His material category can concern a term itself
or another term, including a mental term. His simple category retains
spoken/written reference to concepts. Novaes connects this classification
with Albert's treatment of apparently ambiguous propositions by their
conjoined or disjoined readings, rather than a demand to distinguish
separate senses. The chapter argues for the coherence of this position;
it does not derive that historical interpretation from a modern semantic
theorem. Wyclif and Paul of Venice's realist uses of simple supposition
retain their different ontological basis
\cite[pp.~449--452]{DutilhNovaes2007FordhamChapter8}.

\subsection{Four kinds of descent and the collective case}

Novaes uses the schematic form \(\Phi A\text{ is }\Phi B\), with \(\Phi\)
standing for a quantifying or other syncategorematic expression, or its
absence. Her modern descriptions distinguish:

\begin{enumerate}
\item Propositional conjunctive descent: individual predications are conjoined.
\item Propositional disjunctive descent: individual predications are disjoined.
\item Nominal disjunctive descent: the alternatives remain within a term position.
\item Nominal conjunctive descent: a term position concerns individuals together.
\end{enumerate}

The placement of conjunction and disjunction matters. For every man being
an animal, the subject supports conjunctive individual predications, while
the predicate's alternatives need not pick the same animal for every man.
For no man being an animal, both subject and predicate permit conjunctive
negative descents. For some man being an animal, both permit disjunctive
individual descents. For some man not being an animal, subject descent
is disjunctive; the predicate gives conjunctive negative consequences.
The latter descent does not by itself recover the original shared witness.

\begin{hblemma}{thm:p13-latin-finite-descents}
 {proved=annex:p13.latin.descents}{The witness and scope conditions in finite descents}
In the declared classical model with finite nonempty extensions \(S,P\),
universal and existential individual expansions have their usual conjunctive
and disjunctive forms. From \(\exists x(Sx\land\neg Px)\) one obtains
\(\bigwedge_{b\in P}\exists x(Sx\land x\ne b)\), but the converse fails.
From \(\forall x(Sx\to Px)\) one need not obtain a single \(b\in P\)
equal to every member of \(S\).
\end{hblemma}

These local results explain the scope distinctions without claiming a
complete formalization of every medieval supposition doctrine.
Confused-and-distributive supposition is associated with conjunctive
propositional descent; determinate supposition with disjunctive propositional
descent without the conjunctive form; merely confused supposition with
nominal disjunctive descent not further distinguished by those stronger forms.
The restrictions prevent treating the categories as overlapping merely
because a stronger descent also implies some weaker one.

Maulevelt's example of Socrates differing from every man raises nominal
conjunctive descent. Albert rejects the proposed classification.
Thomas of Cleves associates this descent with collective supposition;
Paul of Venice can instead place it within a broader merely-confused
category. Examples involving not being every man, no animal being every
man, or a penny seen by every man expose scope questions; opponents argue
that other descents still suffice.

The apostles being twelve is a sharper collective example. The group
has the number, not each individual apostle. Proposed responses either
treat the whole collection as a singular object or rephrase the sentence
to force a traditional classification. Ferrer's three common-personal
categories include collective but omit merely confused supposition;
this goes with his restriction of supposition to subjects.
Novaes does not settle whether his classification arose independently
of the Parisian development
\cite[pp.~452--457]{DutilhNovaes2007FordhamChapter8}.

\begin{hblemma}{thm:p13-latin-collective}
 {proved=annex:p13.latin.collective}{A collective predicate is not truth-functional conjunction}
If a predicate of a finite collection means that it has twelve distinct
members, its truth of a twelve-member collection entails neither its truth
of one member nor a disjunction of its truths of the singleton members.
Consequently, an unrestricted nominal-conjunction-to-nominal-disjunction
principle is not justified by the ordinary propositional truth tables.
\end{hblemma}

\subsection{Material occurrences and Marsilius's reordered tree}

Material supposition can concern distinct inscriptions, utterances or
concept acts, not only one word type. Ockham's mature concepts are acts
of conceiving; separate acts and separate inscriptions create a domain
over which one may quantify. Buridan recognizes universal, particular
and singular propositions about such term occurrences.

Marsilius makes the quantitative distinction prior to the personal/material
distinction. His tree begins with discrete versus common supposition;
common divides into determinate and confused; confused divides into
merely confused and confused-and-distributive. Each final quantitative
category has material and personal branches.
For material supposition the examples are a specified red inscription,
at least one occurrence of a word on a page, a nominal disjunction of
word occurrences in an exclusive predication, and a universal predication
of monosyllabicity across those occurrences.
The example that only the word \emph{man} is monosyllabic is a schematic
scope illustration, not a true claim about English vocabulary.
One must distinguish ultimate significates, such as people, from the
non-ultimate significates consisting of terms and term occurrences
\cite[pp.~457--461]{DutilhNovaes2007FordhamChapter8}.

\subsection{Negation requires the whole context}

The ordinary categorical patterns suggest rules for a term following
\emph{every}, \emph{some}, \emph{no}, or a negated copula.
Iterated syncategorematic expressions expose the limitation of a merely
local rule. Retain Novaes's typed notation:
\(\operatorname{Det}(A)_P\), \(\operatorname{Dist}(A)_P\), and
\(\operatorname{Conf}(A)_P\) describe a term's mode in proposition \(P\).
\(\langle\sim,A\rangle_{P^*}\) says that the term follows the added
negation in the modified proposition \(P^*\).
The attributed rule variants are collected in
sheet \ref{sheet:p13-latin-supposition-negation}.

Rules 1 and 2 make formerly determinate or merely-confused terms distributive
when an appropriate negation precedes them. The disputed case is a term
already distributive. Ockham's Rule 3o makes it determinate; Buridan's
Rule 3b makes it merely confused.
Negating no \(A\) being \(B\) should recover some \(A\) being \(B\),
with determinate terms. Negating some \(A\) not being \(B\) should recover
every \(A\) being \(B\), whose predicate is merely confused.
Neither one uniform output for an originally distributive predicate
preserves both desired classifications.

\begin{hblemma}{thm:p13-latin-negation-context}
 {proved=annex:p13.latin.negation-context}{Why one input mode does not determine the negated mode}
If the two displayed contradictory equivalences preserve their intended
supposition classifications and determinate differs from merely confused,
no transformation depending only on the previous distributive mode can
give the required output in both cases.
\end{hblemma}

Dorp's method, reconstructed by Karger from his examples, moves the negation
toward the copula and adjusts the quantifying signs. The Latin example
\emph{Nullum animal omnis homo est} becomes
\emph{Omne animal homo non est}. This is a technical scope transformation,
not permission to read either expression through arbitrary modern English
word order. Novaes's discussion assigns the resulting \emph{animal}
distributive and \emph{homo} merely-confused supposition.
John of St.~Thomas's later, seventeenth-century formulation explicitly
consults the other term's mode: distributive relative to a confused
term becomes determinate under the relevant negation, while distributive
relative to a determinate term becomes merely confused.
This later refinement must not be dated as Ockham's or Buridan's original
rule \cite[pp.~461--465]{DutilhNovaes2007FordhamChapter8}.

\subsection{Proof of a proposition is not automatically a mathematical proof}

The \emph{probationes terminorum} tradition analyzes what a proposition's
truth depends on. Billingham's \emph{Speculum puerorum} is influential,
but is not asserted to be the genre's first work.
Immediate propositions appeal directly to sense or understanding,
often through demonstratives, pronouns and indexicals.
Exponible propositions are analyzed through a conjunction: the example
of Socrates beginning to be white separates earlier non-whiteness from
present whiteness. Resoluble propositions introduce discrete cases:
a running man is addressed through a particular object being a man
and running. Officiable propositions involve accusative-plus-infinitive
nominalizations, including modal and propositional-attitude contexts.
The latter are not licensed by extensional substitution alone.

Novaes explicitly treats the purpose and exact mechanisms of this genre
as incompletely understood. Likewise, the chapter points to further
work on relative-pronoun supposition, dicta and truthmakers rather than
supplying complete systems for them
\cite[pp.~466--467]{DutilhNovaes2007FordhamChapter8}.

\section{Consequences: criteria, rules, and disagreements}
\label{sec:p13-latin-consequences}

\subsection{Origins and four groups of source treatments}

The origin of independent consequence treatises is disputed.
The topical tradition is relevant, but similarity of purpose is not
sufficient evidence that late-thirteenth-century Topics treatises directly
produced the fourteenth-century genre. Novaes considers convergence among
Topics, syncategoremata, hypothetical syllogisms and commentaries on the
\emph{Prior Analytics}. No direct Stoic transmission is established merely
because some rediscovered inference patterns have Stoic counterparts.

The chapter distinguishes four groupings: early collections of rules
associated with Burley and anonymous texts; the more conceptual treatments
using intrinsic/extrinsic middles in Burley, Ockham and related works;
the Buridanian continental substitution criterion; and predominantly British
containment or understanding criteria associated with Fland, Billingham,
Lavenham and Strode. They are broad groupings with exceptions, not an
exhaustive proof of each text's affiliation
\cite[pp.~468--471]{DutilhNovaes2007FordhamChapter8}.

\subsection{Conditional, consequence, and inference}

A conditional is a proposition with embedded clauses. A consequence is
a relation among propositions or assertions. An inference is an action
of drawing a conclusion. These categories can be related without being
identical. Novaes observes that the \emph{Summa}'s consequence treatment
belongs to its argument part, not simply its proposition part.

A widely accepted modal criterion excludes the antecedent's truth together
with the consequent's falsity. Some authors accept it as sufficient;
others impose more. Buridan's actual-proposition commitments motivate a
formulation about what the propositions signify when put forward together.
This is not merely a modern material-implication truth table.
Later writers also allow the name consequence for invalid arguments,
whereas some earlier uses reserve the name for the valid relations.
The source's Pseudo-Scotus self-referential argument about its own invalidity
is an attributed challenge with a cited interpretive dispute, not a
completed general solution of Curry's paradox
\cite[pp.~471--473]{DutilhNovaes2007FordhamChapter8}.

\subsection{The divisions are not one fixed hierarchy}

Ockham's formal/material distinction concerns validating middles.
An extrinsic middle expresses a general connecting rule without the
argument's particular terms; an intrinsic middle uses those terms.
A syllogism can depend on the extrinsic rule alone.
Passing from Socrates not running to a man not running additionally
needs Socrates being a man. The exact text defining Ockham's material
consequences is reported as corrupt, and interpretations restricting
them to impossible antecedents or necessary consequents remain disputed.

Buridan's formal criterion retains the form while varying categorematic
terms, applied within the class already meeting the modal criterion.
Thus a merely accidental current distribution does not establish modal
validity. The example involving a president and being male is a historically
dated illustration of that limitation, not a current demographic claim.
An enthymeme from a man to an animal may receive an understanding-based
formal classification in the British tradition while failing Buridan's
unrestricted term-substitution test until its missing universal premise
is included.

Paul of Venice uses a conceivability condition and also distinguishes
\emph{bona de forma} from \emph{bona de materia}. Novaes records tensions
in how the labels are applied, including the running-man/running-animal
example; terminology alone does not erase those tensions.
Absolute versus as-of-now consequence concerns the stability of the
supporting conditions. Pseudo-Scotus applies it to material consequences
through necessary versus contingent missing premises; Peter of Mantua
orders the distinctions differently and uses the time designated by the
copula. These are overlapping classifications, not one silently unified
tree \cite[pp.~474--478]{DutilhNovaes2007FordhamChapter8}.

\subsection{The source's consequence operations and full local derivations}

Novaes introduces \(\Rightarrow\) as a modern representation of consequence,
not as the conditional connective \(\to\).
Her rule displays are collected in
sheet \ref{sheet:p13-latin-consequence-operations}.
For the local proofs, fix a class \(W\) of classical valuations.
\(A\Rightarrow B\) means that every \(v\in W\) making \(A\) true makes
\(B\) true. Multiple antecedents mean simultaneous truth; multiple
consequents in the source's derived-rule discussion are joint conclusions.
This latter comma is not the disjunctive right comma of classical LK.
\(T\) denotes a proposition true throughout \(W\), and \(\bot\) one
false throughout \(W\).
Restricting \(W\) to a specified temporal/background case gives a
qualified as-of-now reconstruction; it is not a theorem about all possible
backgrounds.

\begin{hblemma}{thm:p13-latin-consequence-operations}
 {proved=annex:p13.latin.consequence-operations}{The displayed consequence operations in their declared model}
In this model, transitivity, both displayed derived cuts, concession and
denial propagation, contraposition and its two composite forms,
conjunction projection, disjunction introduction, and conditional detachment
hold. Impossible antecedents and necessary consequents give the two
vacuous boundary cases. Both corresponding non-consequence
propagation principles have countermodel proofs.
\end{hblemma}

The first derived cut uses \(A\Rightarrow B\) to obtain joint
\(A,B\) from \(A\), then applies the premise \(A,B\Rightarrow C\).
The second adds the same \(C\) to the antecedent and joint conclusion
before applying \(B,C\Rightarrow D\).
The source calls these intermediate steps special cases of weakening;
the proof below uses their actual joint-conclusion meaning.
The two non-consequence principles transport different countervaluations:
one changes the antecedent along a valid consequence, and the other changes
the failed consequent. Each has its own explicit countermodel argument.

Term relations also matter. If men are included among animals, the
universal predication that every animal runs entails that every man runs.
Existence of a running man additionally needs a nonempty class of men.
The reverse inference does not hold merely because every man runs.

\begin{hblemma}{thm:p13-latin-superior}
 {proved=annex:p13.latin.superior}{Superior and inferior terms with existential import}
If \(M\subseteq A\) and \(A\subseteq R\), then \(M\subseteq R\).
If also \(M\ne\varnothing\), some member of \(M\) belongs to \(R\).
Neither \(M\subseteq R\) nor the existence of such a member entails
\(A\subseteq R\).
\end{hblemma}

\subsection{Dissident restrictions are substantive}

Novaes reports restrictions by Nicolaus Drukken of Dacia and Richard
Ferrybridge, and antecedents in Abelard and Kilwardby.
Containment of what is signified, or a required relevance relation,
can block unrelated conclusions from an impossible antecedent or a
necessary consequent. This rejects sufficiency of the unrestricted
modal criterion; it is not a rejection of every inferential relation.
But merely announcing relevance while retaining all explosion-producing
rules does not change their deductive effect.

\begin{hblemma}{thm:p13-latin-explosion}
 {proved=annex:p13.latin.explosion}{The complete explosion-producing chain}
Conjunction projection, disjunction introduction and disjunctive syllogism
together derive arbitrary \(Q\) from \(P\land\neg P\).
Any system refusing that derivation must restrict at least one of those
steps, their applicable contexts, or their composition.
\end{hblemma}

The chapter does not give a complete calculus for the dissident theories.
The local derivation locates the dependency; it does not manufacture
an ancient paraconsistent calculus
\cite[pp.~479--484]{DutilhNovaes2007FordhamChapter8}.

\section{Obligations: Burley's expanding commitments}
\label{sec:p13-latin-burley-obligations}

\subsection{What the respondent undertakes}

An obligation involves an Opponent proposing statements and a Respondent
conceding, denying or doubting under stipulated rules. Novaes traces a
thirteenth-century Parisian background and a predominantly British
fourteenth-century development; the genre is not invented by its first
surviving fourteenth-century representative.
Burley lists \emph{petitio}, \emph{sit verum}, \emph{institutio},
\emph{positio}, \emph{depositio}, and \emph{dubitatio}.
Swyneshed lists \emph{positio}, \emph{depositio}, and \emph{impositio}.
The detailed constructions here concern \emph{positio}, not all these forms.

In the chapter's chronology, Burley's obligation treatise belongs to the
century's opening years; Kilvington's pre-1325 discussion precedes
Swyneshed's approximately 1330--1335 innovation. Fland uses the older/newer
response labels around 1350. Billingham, Martin of Alnwick and Brinkley
are placed in the 1350s and Strode in the 1360s.
Martin, Lavenham and to some extent Fland follow the new response style;
the other listed British treatments broadly retain the old style.
For the continent, the account places Albert and John of Wesel around 1350,
Buser around 1355, Marsilius just before 1360, Holland just after 1360,
and Peter of Candia at the century's end. The last two had especially
direct British exposure; the others' dependence on British discussions
does not prove agreement with Swyneshed. Most criticize his relevance
change. These are the chapter's qualified dating and reception claims.

In Novaes's Burley reconstruction,
\[
 \operatorname{Ob}=\langle K_C,\Phi,\Gamma,R(\phi_n)\rangle .
\]
\(K_C\) is fixed shared information together with an assumed \emph{casus}.
\(\Phi\) records the ordered proposals, beginning with the positum
\(\phi_0\). Repetitions are occurrences in this history, not automatically
distinct propositions. \(\Gamma_n\) is the set of commitments after round
\(n\). A concession adds \(\phi_n\), a denial adds \(\neg\phi_n\),
and doubt adds neither. \(R\) has values \(1,0,?\).
The source uses \(\Vdash\) for semantic following, not set-theoretic forcing.
The function notation depends implicitly on the current round and history.
Here \(\phi\Vdash\psi\) abbreviates following from the singleton
premise set \(\{\phi\}\); the source glyph does not turn an object-language
formula into a premise set without this convention.

The positum is admitted only when it is not contradictory.
For later proposals, inferential relations with \(\Gamma_{n-1}\) take
priority. If the proposal follows, concede it; if its contradictory
follows, deny it. If neither follows, answer from \(K_C\): concede what
is known true, deny what is known false, and doubt what is unsettled.
Rules are given in sheet \ref{sheet:p13-latin-burley-obligations}.
Known information must itself not force both sides of the same question.
The positum may be false according to \(K_C\); the two bases are not
simply united. The source example supposes Socrates black while positing
him white \cite[pp.~484--489]{DutilhNovaes2007FordhamChapter8}.

\begin{hblemma}{thm:p13-latin-burley-consistency}
 {proved=annex:p13.latin.burley-consistency}{Correct responses preserve the admitted commitment base}
In classical semantics, with a satisfiable positum and consistent shared
information, the stated response conditions give a unique response at
each reached state and preserve satisfiability of \(\Gamma_n\).
This supplies a logical winning response policy for every finite sequence
of correctly evaluated proposals, not an efficient human strategy.
\end{hblemma}

The full consistent-extension argument matters: if neither side is already
forced by the current base, either side can be added consistently.
This is the Lindenbaum-style fact identified in the source's footnote 65.
Unlike an unchecked maximal-set construction, an actual wrong commitment
can lose a disputation; the rule does not erase it after the event.
The game is time-limited and need not construct a maximal theory.

\subsection{The Rome example: a hidden inferential dependency}

Novaes's worked trap is explicitly her illustration, not a quotation of
a particular Burley disputation. Use \(a\) for the respondent, \(b\)
for the pope, and \(L(x)\) for being in Rome.
Shared information includes \(\neg L(a)\) and \(L(b)\).
The proposed sequence and the tempting responses are:

\begin{enumerate}
\item \(L(a)\lor a=b\): concede as a consistent positum.
\item \(L(a)\): deny as initially independent and known false.
\item \(L(b)\): mistakenly concede as independent and known true.
\item \(a=b\): concede because it follows from the positum and the denial.
\item \(L(a)\): concede from equality and \(L(b)\), displaying the contradiction.
\end{enumerate}

\begin{hblemma}{thm:p13-latin-rome}
 {proved=annex:p13.latin.rome}{The complete dependency that defeats the apparent trap}
After the second response, the commitment base already entails
\(a=b\) and \(\neg L(b)\). Therefore the third proposal must be denied,
not answered from shared information. The correct five-round trace
remains satisfiable; the tempting trace does not.
\end{hblemma}

This example presupposes the equality/substitution relations in its
language. Treating its three verbal statements as unrelated Boolean
atoms would omit the very dependency that makes the trap work.
Being in Rome is a stipulated historical-case predicate here, not a
claim about a current office holder.

\subsection{Models and expanding information}

Retain the source notation \(\mathrm{UM}_n\) for all models satisfying
\(\Gamma_n\), and \(\mathrm{UM}_{\phi}\) for all models satisfying \(\phi\).
Adding a literal \(\psi\) gives
\[
 \mathrm{UM}_{n+1}=\mathrm{UM}_n\cap\mathrm{UM}_{\psi}.
\]

\begin{hblemma}{thm:p13-latin-model-restriction}
 {proved=annex:p13.latin.model-restriction}{Both source model-restriction results}
If the added literal is already entailed, then
\(\mathrm{UM}_{n+1}=\mathrm{UM}_n\).
If a previously independent proposal is conceded or denied, then
\(\varnothing\ne\mathrm{UM}_{n+1}\subsetneq\mathrm{UM}_n\).
Doubt leaves the model class unchanged.
\end{hblemma}

These are the source's Theorems 9 and 10, with the deny/doubt cases made
explicit. Commitments grow while their permissible models shrink.
Later relevance can therefore depend on earlier proposals and their order.
For instance, from a positum \(p\lor q\), conceding independent \(p\)
need not settle \(q\), while denying independent \(p\) forces \(q\).
The inference does not reverse: positing \(p\lor q\) and conceding \(q\)
does not by itself force \(\neg p\).

\section{Swyneshed: fixed inferential relevance and variable knowledge}
\label{sec:p13-latin-swyneshed-obligations}

\subsection{A different target of the exercise}

Novaes calls Burley's pattern \emph{antiqua responsio} and Swyneshed's
\emph{nova responsio}, while warning that historical uses of older/newer
labels vary by generation. Swyneshed's changes concern relevance to the
positum alone rather than to all prior concessions and denials.
The chapter relates this to inference recognition rather than maintenance
of the union of all responses.

His reconstructed data are
\[
 \operatorname{Ob}=\langle\Sigma,\Phi,I,R(\phi_n)\rangle .
\]
\(\Sigma\) is a sequence of knowledge states \(S_n\); \(I\) is the
ordered reply history with records \(\iota_n=[\phi;\gamma]\),
\(\gamma\in\{1,0,?\}\). The proposition identifier can recur at different
response indices. The reply history is not the expanding inferential base.
The initial positum must be contingent, excluding both impossibility and
necessity. This prevents a trivial inference-recognition exercise.

The chapter also proposes a reply-history formulation: reject a candidate
whose replies are uniformly one \(\gamma\), and admit one for which
there are replies \(\gamma\ne\kappa\) at moments \(n,m\).
This description must not be used as a finite effective test of modal
contingency. A contingent proposition can remain known true throughout
an actual finite record, and a repeatedly proposed admitted positum
normally receives the same required response. The operative admissibility
condition below is the stated modal contingency, not an inference
from absence of observed reply variation.

The source's proposed history conditions can be recorded, without
using them as an admission algorithm, as
\[
 R(\phi_0)=0\quad\text{iff all compared replies }
 \iota_n=[\phi_0;\gamma],\ \iota_m=[\phi_0;\gamma]
 \text{ have one }\gamma,
\]
\[
 R(\phi_0)=1\quad\text{iff some compared replies }
 \iota_n=[\phi_0;\gamma],\ \iota_m=[\phi_0;\kappa]
 \text{ have }\gamma\ne\kappa .
\]
The modal counterexample in the annex separates the intended contingency
criterion from this reading of a uniform finite reply history.

With a fixed proposition universe \(\mathcal L\), define
\[
 \begin{aligned}
 \Delta_{\phi_0}&=
 \{\phi\in\mathcal L:\phi_0\Vdash\phi\text{ or }\phi_0\Vdash\neg\phi\},
 \\
 \Pi_{\phi_0}&=
 \{\phi\in\mathcal L:\phi_0\nVdash\phi\text{ and }\phi_0\nVdash\neg\phi\}.
 \end{aligned}
\]
For pertinent proposals, concede consequences of the positum and deny
their repugnants. For impertinent proposals, consult the present
knowledge state \(S_n\), including its possible ignorance.
The sheet is \ref{sheet:p13-latin-swyneshed-obligations}.
No prior impertinent reply is added to the pertinence test
\cite[pp.~492--496]{DutilhNovaes2007FordhamChapter8}.

The ordinary procedural alternation remains proposal followed by reply.
Opponent ends the exercise by declaring \emph{Cedat tempus obligationis}.
A wrongly evaluated reply can lose the exercise; otherwise the final
assessment can recognize correct performance. The criterion is not
consistency of the union of every logged response, as the next examples show.

\begin{hblemma}{thm:p13-latin-swyneshed-response}
 {proved=annex:p13.latin.swyneshed-response}{Response determination with its state qualification}
For a satisfiable positum and consistent current knowledge state, the two
classes are disjoint and exhaustive and the rule assigns one reply.
For fixed positum, proposal and knowledge state, changing the preceding
reply history does not change that reply. Changing \(S_n\) can change
an impertinent reply.
\end{hblemma}

The exceptional positum that nothing is posited to the respondent is
pragmatically inconsistent with the act of positing. Swyneshed treats
its repetition as impertinent and denies it; Burley would reject such
a positum initially according to the chapter's comparison.
This is a special pragmatics-of-disputation provision, not a contradiction
in the ordinary fixed-\(\Delta\) definition.
Adding a second positum likewise changes the base for pertinence;
the stated single-positum model does not pretend otherwise.

\subsection{Why concessions do not form one truth assignment}

A correct new-style record can grant \(p\) and \(q\) but deny \(p\land q\),
or grant \(p\lor q\) while denying each disjunct.
The statements have different sources of entitlement: some are answered
relative to the positum, others relative to what is currently known.

\begin{hblemma}{thm:p13-latin-swyneshed-connectives}
 {proved=annex:p13.latin.swyneshed-connectives}{Complete contrasting reply records}
With positum \(\neg(p\land q)\) and current knowledge \(p\land q\),
the correct replies to \(p,q,p\land q\) are \(1,1,0\).
With positum \(p\lor q\) and current knowledge \(\neg p\land\neg q\),
the correct replies to \(p,q,p\lor q\) are \(0,0,1\).
The corresponding unqualified unions of commitments are inconsistent.
Neither trace invalidates the classical truth conditions within either
semantic evaluation.
\end{hblemma}

The source's two-column image separates pertinent and impertinent responses.
Conjunction closure holds where the relevant evaluations and classifications
permit it. But imperfect information adds a further limit to the claim
that granting a disjunction requires granting a disjunct.

\begin{hblemma}{thm:p13-latin-partial-knowledge}
 {proved=annex:p13.latin.partial-knowledge}{A known disjunction need not have a known disjunct}
In the fixed knowledge state consisting of consequences of \(p\lor q\),
neither \(p\) nor \(q\) is known individually. With independent positum
\(r\), all three formulas are impertinent: the replies are \(1,?,?\)
to \(p\lor q,p,q\). Thus the disjunct-choice assertion requires more
than consistent but incomplete knowledge.
\end{hblemma}

This limit concerns the modern knowledge-state reconstruction, not a
claim that every historical knowledge judgment is classical entailment.
Changes in things can also produce opposite replies at separate times:
the respondent may first be seated and later standing. Swyneshed consults
the relevant current state for each impertinent proposal; Burley's earlier
commitment takes priority in his own counterfactual disputation.
The temporal or inferential indexing must therefore be retained before
calling a combined response record contradictory.

Novaes's conclusion is comparative, not a complete account of every
obligation genre or sophism. The chapter leaves further self-reference
and pragmatic puzzles to its cited literature. Its local semantic
reconstructions illuminate differences in reasoning procedures without
identifying commitment, knowledge and historical interpretation.

\subsection{Supposition, consequence, and obligation rule sheets}
% Novaes, Fordham manuscript chapter8: modern reconstructions, not ancient notation.
\calculussheet{sheet:p13-latin-supposition-negation}
 {Novaes's attributed supposition-negation variants}
 {A and B are terms; P is their original proposition and P-star its negated modification.
 Det, Dist and Conf classify a term in a whole proposition. These variants are not one calculus.}
\hbrule{rule:p13-latin-neg-1}{Rule 1}
 {\operatorname{Det}(A)_P\quad\langle\sim,A\rangle_{P^*}}
 {\operatorname{Dist}(A)_{P^*}}{\(A\) follows the relevant added negation.}
\hbrule{rule:p13-latin-neg-2}{Rule 2}
 {\operatorname{Conf}(A)_P\quad\langle\sim,A\rangle_{P^*}}
 {\operatorname{Dist}(A)_{P^*}}{\(A\) follows the relevant added negation.}
\hbrule{rule:p13-latin-neg-3o}{Rule 3o: Ockham variant}
 {\operatorname{Dist}(A)_P\quad\langle\sim,A\rangle_{P^*}}
 {\operatorname{Det}(A)_{P^*}}{Attributed variant, not combined with Rule 3b.}
\hbrule{rule:p13-latin-neg-3b}{Rule 3b: Buridan variant}
 {\operatorname{Dist}(A)_P\quad\langle\sim,A\rangle_{P^*}}
 {\operatorname{Conf}(A)_{P^*}}{Attributed variant, not combined with Rule 3o.}
\hbrule{rule:p13-latin-neg-3op}{Contextual Rule 3o-prime}
 {\operatorname{Dist}(A)_P\quad\operatorname{Conf}(B)_P\quad
  \langle\sim,A\rangle_{P^*}}
 {\operatorname{Det}(A)_{P^*}}
 {\(B\) is the other term; John of St. Thomas's later refinement.}
\hbrule{rule:p13-latin-neg-3bp}{Contextual Rule 3b-prime}
 {\operatorname{Dist}(A)_P\quad\operatorname{Det}(B)_P\quad
  \langle\sim,A\rangle_{P^*}}
 {\operatorname{Conf}(A)_{P^*}}
 {\(B\) is the other term; John of St. Thomas's later refinement.}

\calculussheet{sheet:p13-latin-consequence-operations}
 {Novaes's displayed consequence operations and declared local semantics}
 {Formula letters A B C D; consequence arrow distinct from the conditional connective.
 Fixed class W of classical valuations; left comma simultaneous premises,
 right comma joint conclusions in the source's derived-cut discussion, not LK right comma.}
\hbrule{rule:p13-latin-cut}{Burley Rule 2}
 {A\Rightarrow B\quad B\Rightarrow C}{A\Rightarrow C}{The same class \(W\) throughout.}
\hbrule{rule:p13-latin-cut-prime}{Rule 2-prime}
 {A\Rightarrow B\quad A,B\Rightarrow C}{A\Rightarrow C}{Joint premises in the second consequence.}
\hbrule{rule:p13-latin-cut-double-prime}{Rule 2-double-prime}
 {A\Rightarrow B\quad B,C\Rightarrow D}{A,C\Rightarrow D}{The added \(C\) is retained.}
\hbrule{rule:p13-latin-joint-conclusion}{Source WR step}
 {A\Rightarrow B}{A\Rightarrow A,B}
 {The right side means both conclusions; this is not unrestricted LK weakening.}
\hbrule{rule:p13-latin-add-same}{Source WR/WL step}
 {A\Rightarrow B}{A,C\Rightarrow B,C}{The same \(C\) on both sides.}
\hbrule{rule:p13-latin-impossible}{Ex impossibili}
 {}{\bot\Rightarrow A}{\(\bot\) has no verifying valuation in \(W\).}
\hbrule{rule:p13-latin-necessary}{Ad necessarium}
 {}{A\Rightarrow T}{\(T\) is true throughout \(W\).}
\hbrule{rule:p13-latin-concession}{Concession propagation}
 {\Rightarrow A\quad A\Rightarrow B}{\Rightarrow B}{Empty antecedent means truth throughout \(W\).}
\hbrule{rule:p13-latin-denial}{Denial propagation}
 {\Rightarrow\neg B\quad A\Rightarrow B}{\Rightarrow\neg A}{Classical contradictories.}
\hbrule{rule:p13-latin-contraposition}{Contraposition}
 {A\Rightarrow B}{\neg B\Rightarrow\neg A}{Classical contradictories.}
\hbrule{rule:p13-latin-negative-chain}{Negative chain}
 {A\Rightarrow B\quad\neg A\Rightarrow C}{\neg B\Rightarrow C}{The same consequence criterion.}
\hbrule{rule:p13-latin-negative-predecessor}{Negative predecessor}
 {A\Rightarrow B\quad C\Rightarrow\neg B}{C\Rightarrow\neg A}{The same consequence criterion.}
\hbrule{rule:p13-latin-conjunction-projection}{Conjunction projection}
 {}{A\mathbin{\&}B\Rightarrow A}{The second projection follows by exchange of letters.}
\hbrule{rule:p13-latin-disjunction-introduction}{Disjunction introduction}
 {}{A\Rightarrow A\lor B}{The displayed source pattern.}
\hbrule{rule:p13-latin-conditional-detachment}{Conditional detachment}
 {}{A\to B,A\Rightarrow B}{The object connective uses the declared classical local reading.}
\hbrule{rule:p13-latin-negative-successor}{Non-consequence: successor}
 {A\Rightarrow B\quad A\nRightarrow C}{B\nRightarrow C}
 {The same class \(W\); the non-consequence premise supplies an explicit countervaluation.}
\hbrule{rule:p13-latin-negative-antecedent}{Non-consequence: antecedent}
 {A\Rightarrow B\quad C\nRightarrow B}{C\nRightarrow A}
 {Explicit countervaluation witnesses non-consequence.}
\hbrule{rule:p13-latin-superior-universal}{Superior to inferior: universal}
 {M\subseteq A\quad A\subseteq R}{M\subseteq R}{Declared extension interpretation of the term relation.}
\hbrule{rule:p13-latin-superior-existential}{Superior to inferior: existential}
 {M\subseteq A\quad A\subseteq R\quad M\ne\varnothing}
 {\exists x(x\in M\land x\in R)}{Existential import is explicit.}
\hbrule{rule:p13-latin-disjunctive-syllogism}{Explosion-chain dependency}
 {}{P\lor Q,\neg P\Rightarrow Q}
 {The source's Lewis-style argument uses this step; dissident scope restrictions remain separate.}

\calculussheet{sheet:p13-latin-burley-obligations}
 {Novaes's Burley-style positio reconstruction}
 {Semantic following uses the source Vdash glyph, not set-theoretic forcing.
 Gamma is the current commitment base; KC is separate consistent static shared information;
 response values 1 0 ? mean concede deny doubt.}
\hbrule{rule:p13-latin-burley-initial-yes}{Admit positum}
 {\phi_0\nVdash\bot}{R(\phi_0)=1}{Initialize \(\Gamma_0=\{\phi_0\}\).}
\hbrule{rule:p13-latin-burley-initial-no}{Reject contradictory positum}
 {\phi_0\Vdash\bot}{R(\phi_0)=0}{No disputation starts.}
\hbrule{rule:p13-latin-burley-forced-yes}{Pertinent concession}
 {\Gamma_{n-1}\Vdash\phi_n}{R(\phi_n)=1}
 {\(n>0\); add \(\phi_n\) to the base.}
\hbrule{rule:p13-latin-burley-forced-no}{Pertinent denial}
 {\Gamma_{n-1}\Vdash\neg\phi_n}{R(\phi_n)=0}
 {\(n>0\); add \(\neg\phi_n\) to the base.}
\hbrule{rule:p13-latin-burley-known-yes}{Impertinent known true}
 {\Gamma_{n-1}\nVdash\phi_n\quad\Gamma_{n-1}\nVdash\neg\phi_n\quad
  K_C\Vdash\phi_n}{R(\phi_n)=1}{Add only \(\phi_n\), not all of \(K_C\).}
\hbrule{rule:p13-latin-burley-known-no}{Impertinent known false}
 {\Gamma_{n-1}\nVdash\phi_n\quad\Gamma_{n-1}\nVdash\neg\phi_n\quad
  K_C\Vdash\neg\phi_n}{R(\phi_n)=0}{Add only \(\neg\phi_n\).}
\hbrule{rule:p13-latin-burley-unknown}{Impertinent unknown}
 {\Gamma_{n-1}\nVdash\phi_n\quad\Gamma_{n-1}\nVdash\neg\phi_n\quad
  K_C\nVdash\phi_n\quad K_C\nVdash\neg\phi_n}
 {R(\phi_n)=?}{Leave \(\Gamma_n=\Gamma_{n-1}\).}

\calculussheet{sheet:p13-latin-swyneshed-obligations}
 {Novaes's Swyneshed-style positio reconstruction}
 {Fixed single positum; Delta and Pi determine semantic pertinence.
 Sn is current consistent, possibly incomplete information. The response log is not an inferential base.
 Pragmatically repugnant re-positing is an explicitly separate exception.}
\hbrule{rule:p13-latin-swyneshed-initial-yes}{Contingent positum}
 {\phi_0\nVdash\bot\quad\neg\phi_0\nVdash\bot}{R(\phi_0)=1}
 {The modal condition, not a finite reply-history variation test.}
\hbrule{rule:p13-latin-swyneshed-initial-no}{Noncontingent candidate}
 {\phi_0\Vdash\bot\ \lor\ \neg\phi_0\Vdash\bot}{R(\phi_0)=0}
 {Exclude impossible and necessary candidates in the declared classical modal class.}
\hbrule{rule:p13-latin-swyneshed-repeat}{Ordinary repeated positum}
 {\phi_n=\phi_0}{R(\phi_n)=1}{No pragmatic re-positing exception applies.}
\hbrule{rule:p13-latin-swyneshed-pertinent-yes}{Pertinent concession}
 {\phi_n\in\Delta_{\phi_0}\quad\phi_0\Vdash\phi_n}{R(\phi_n)=1}
 {\(\phi_n\ne\phi_0\); the base is the positum alone.}
\hbrule{rule:p13-latin-swyneshed-pertinent-no}{Pertinent denial}
 {\phi_n\in\Delta_{\phi_0}\quad\phi_0\Vdash\neg\phi_n}{R(\phi_n)=0}
 {\(\phi_n\ne\phi_0\); the base is the positum alone.}
\hbrule{rule:p13-latin-swyneshed-known-yes}{Impertinent known true}
 {\phi_n\in\Pi_{\phi_0}\quad S_n\Vdash\phi_n}{R(\phi_n)=1}
 {Record the reply; do not enlarge the pertinence base.}
\hbrule{rule:p13-latin-swyneshed-known-no}{Impertinent known false}
 {\phi_n\in\Pi_{\phi_0}\quad S_n\Vdash\neg\phi_n}{R(\phi_n)=0}
 {Record the reply; do not enlarge the pertinence base.}
\hbrule{rule:p13-latin-swyneshed-unknown}{Impertinent unknown}
 {\phi_n\in\Pi_{\phi_0}\quad S_n\nVdash\phi_n\quad S_n\nVdash\neg\phi_n}
 {R(\phi_n)=?}{Imperfect information remains allowed.}
\hbrule{rule:p13-latin-swyneshed-pragmatic-repeat}{Pragmatically repugnant re-positing}
 {\phi_n=\phi_0\quad\phi_n\text{ is repugnant to the act of positing}}
 {R(\phi_n)\text{ is evaluated from }S_n}
 {Exceptional historical provision, not ordinary membership in the fixed Delta partition.}

